\documentclass[12pt]{article}
\usepackage[utf8]{inputenc}
\usepackage[T1]{fontenc}
\usepackage{hyperref}
\title{Kabatei Language Resources for K-12 Curriculum Development}
\author{Compiled by Claude Sonnet 5.5 (claude-sonnet-5-5) for Andr\'as Kornai}
\date{}
\begin{document}
\maketitle

\section{Introduction}
Kabatei (ISO 639-3 xkp; Glottocode kaba1276) is an Indo-European language of the Tatic group of Northwestern Iranian. Glottolog places it in Taromic, and Ethnologue-derived data locates it in Rudbar District, Gilan Province, Iran. Dialects listed are Kalas and Kabate. It is also called Kafteji, and it is close to Kelasi and Upper Taromi. Verbix lists its vulnerability as Dormant with no speaker estimate, which is more pessimistic than the survey's description (first language of adults only).

\section{Primary Educational Resources}
No primers, textbooks or school programmes for Kabatei were found in the searches.

\section{Linguistic and Reference Materials}
\begin{itemize}
\item Glottolog, languoid kaba1276 (classification and bibliography entry point). \url{https://glottolog.org/resource/languoid/id/kaba1276.bigmap.html}
\item WALS, entry for Kabatei (code kbt): location and typological data. \url{https://wals.info/languoid/lect/wals_code_kbt}
\item Wikipedia, ``Tati language (Iran)'': overview of the Tati umbrella, including xkp. \url{https://en.wikipedia.org/wiki/Tati_language_(Iran)}
\item Verbix, Endangered Languages: Kabatei. \url{https://docs.verbix.com/EndangeredLanguages/Kabatei}
\item Endangered Languages Project record. \url{https://www.endangeredlanguages.com/lang/3247?hl=de}
\item ``Essays on typology of Iranian Languages'' (library record): a Comrie-edited volume that discusses Kafteji (Kabatei) and its gender system. \url{https://vufind.lib.tsu.ru/Record/libtsu.1013260/Details}
\end{itemize}

\section{Religious and Community Resources}
No religious or community-organisation resources were found.

\section{Digital and Audio Resources}
No audio, video or corpus resources were found in the search results. The Endangered Languages Project record may link to samples, but this was not checked.

\section{Curriculum Development Context}
Documentation is thin and the language is possibly dormant. Fieldwork with remaining speakers would come first. Any school use in Iran would also need to be arranged within the national education framework. The names Kabatei, Kafteji and Kalas/Kelasi overlap in sources, so these varieties must be kept apart.

\end{document}